%\documentclass[tikz,border=10pt]{standalone}
%\usepackage{amsmath}
%\usetikzlibrary{positioning,arrows.meta,fit,calc,backgrounds}
%
%\begin{document}
%	\begin{tikzpicture}[
%		font=\small,
%		>=Latex,
%		box/.style={
%			draw,
%			rounded corners=3pt,
%			thick,
%			align=left,
%			inner sep=6pt,
%			text width=5.2cm,
%			minimum height=1.0cm
%		},
%		titlebox/.style={
%			draw,
%			rounded corners=3pt,
%			very thick,
%			align=center,
%			inner sep=6pt,
%			text width=5.4cm,
%			minimum height=0.9cm,
%			font=\bfseries
%		},
%		note/.style={
%			draw,
%			rounded corners=3pt,
%			dashed,
%			align=left,
%			inner sep=5pt,
%			text width=5.2cm
%		},
%		flow/.style={->, thick},
%		compare/.style={->, very thick, dashed}
%		]
%		
%		% ------------------------------------------------------------------
%		% Left column: Symbolic security
%		% ------------------------------------------------------------------
%		\node[titlebox] (symtitle) at (0,0) {Symbolic Security};
%		
%		\node[box, below=0.5cm of symtitle] (symmsg) {
%			\textbf{Messages}\\
%			Represented as \emph{symbolic terms}
%		};
%		
%		\node[box, below=0.35cm of symmsg] (symprim) {
%			\textbf{Cryptographic primitives}\\
%			Idealized as black-box functions with equations
%		};
%		
%		\node[box, below=0.35cm of symprim] (symeq) {
%			\textbf{Canonical example}\\[2pt]
%			\[
%			\mathrm{Dec}(\mathrm{Enc}(m,k),k)=m
%			\]
%		};
%		
%		\node[box, below=0.35cm of symeq] (symadv) {
%			\textbf{Adversary model}\\
%			Manipulates terms using symbolic rules and the available equations
%		};
%		
%		\node[note, below=0.45cm of symadv] (symgain) {
%			\textbf{Main gain}\\
%			Tractable automated reasoning about:
%			\begin{itemize}
%				\item protocol logic
%				\item message flows
%				\item authentication failures
%				\item secrecy violations
%				\item state-machine flaws
%			\end{itemize}
%		};
%		
%		% ------------------------------------------------------------------
%		% Right column: Computational security
%		% ------------------------------------------------------------------
%		\node[titlebox] (comptitle) at (8.6,0) {Computational Security};
%		
%		\node[box, below=0.5cm of comptitle] (compmsg) {
%			\textbf{Messages}\\
%			Represented as \emph{bitstrings}
%		};
%		
%		\node[box, below=0.35cm of compmsg] (compprim) {
%			\textbf{Cryptographic primitives}\\
%			Specified as probabilistic algorithms
%		};
%		
%		\node[box, below=0.35cm of compprim] (compalg) {
%			\textbf{Canonical example}\\[2pt]
%			\[
%			(\mathsf{Gen},\mathsf{Enc},\mathsf{Dec})
%			\]
%			Security is defined through formal experiments
%		};
%		
%		\node[box, below=0.35cm of compalg] (compadv) {
%			\textbf{Adversary model}\\
%			Probabilistic polynomial-time attacker\\
%			(or concrete attacker with explicit resource bounds)
%		};
%		
%		\node[note, below=0.45cm of compadv] (compgain) {
%			\textbf{Main gain}\\
%			Faithful analysis in terms of:
%			\begin{itemize}
%				\item probabilities
%				\item reductions
%				\item adversarial resources
%				\item hardness assumptions
%			\end{itemize}
%		};
%		
%		% ------------------------------------------------------------------
%		% Vertical flows
%		% ------------------------------------------------------------------
%		\draw[flow] (symtitle) -- (symmsg);
%		\draw[flow] (symmsg) -- (symprim);
%		\draw[flow] (symprim) -- (symeq);
%		\draw[flow] (symeq) -- (symadv);
%		\draw[flow] (symadv) -- (symgain);
%		
%		\draw[flow] (comptitle) -- (compmsg);
%		\draw[flow] (compmsg) -- (compprim);
%		\draw[flow] (compprim) -- (compalg);
%		\draw[flow] (compalg) -- (compadv);
%		\draw[flow] (compadv) -- (compgain);
%		
%		% ------------------------------------------------------------------
%		% Bottom comparison
%		% ------------------------------------------------------------------
%		\node[box, text width=5.0cm, minimum height=1.3cm, below=1.0cm of symgain] (symtrade) {
%			\centering
%			\textbf{Strength}\\
%			More abstract, more tractable, easier to automate
%		};
%		
%		\node[box, text width=5.0cm, minimum height=1.3cm, below=1.0cm of compgain] (comptrade) {
%			\centering
%			\textbf{Strength}\\
%			More realistic, closer to standard provable cryptography
%		};
%		
%		\draw[flow] (symgain) -- (symtrade);
%		\draw[flow] (compgain) -- (comptrade);
%		
%		\draw[compare] (symtrade.east) -- node[above, align=center] {\textbf{Trade-off}} (comptrade.west);
%		
%		% ------------------------------------------------------------------
%		% Overall caption-like label
%		% ------------------------------------------------------------------
%		\node[draw, rounded corners=3pt, very thick, align=center,
%		inner sep=6pt, text width=12.2cm, below=1.1cm of $(symtrade)!0.5!(comptrade)$] (summary) {
%			\textbf{Core contrast}\\
%			Symbolic security abstracts cryptography into terms and equations to enable strong automation;\\
%			computational security models algorithms, probabilities, and resources more faithfully, but is generally harder to automate completely.
%		};
%		
%		% Background grouping
%		\begin{scope}[on background layer]
%			\node[draw, rounded corners=6pt, line width=0.9pt,
%			fit=(symtitle)(symtrade), inner sep=8pt] {};
%			\node[draw, rounded corners=6pt, line width=0.9pt,
%			fit=(comptitle)(comptrade), inner sep=8pt] {};
%		\end{scope}
%		
%	\end{tikzpicture}
%\end{document}

\documentclass[tikz,border=10pt]{standalone}
\usepackage{amsmath,amssymb,mathtools}
\usetikzlibrary{positioning,arrows.meta,fit,calc,backgrounds}

\begin{document}
	
	\begin{tikzpicture}[
		>=Latex,
		font=\small,
		paneltitle/.style={
			draw,
			rounded corners=3pt,
			very thick,
			align=center,
			inner sep=6pt,
			minimum height=10mm,
			text width=7.0cm,
			font=\bfseries
		},
		block/.style={
			draw,
			rounded corners=3pt,
			thick,
			align=left,
			inner sep=6pt,
			text width=7.0cm
		},
		smallblock/.style={
			draw,
			rounded corners=3pt,
			thick,
			align=left,
			inner sep=5pt,
			text width=6.7cm
		},
		arrow/.style={->, thick},
		compare/.style={->, very thick, dashed}
		]
		
		% =========================================================
		% LEFT: SYMBOLIC
		% =========================================================
		\node[paneltitle] (symtitle) at (0,0) {Symbolic Security};
		
		\node[block, below=6mm of symtitle] (sym1) {%
			\textbf{Message universe as a term algebra.}\\[2pt]
			Let $\Sigma$ be a cryptographic signature, e.g.
			\[
			\Sigma \supseteq \{\mathsf{Enc},\mathsf{Dec},\mathsf{pair},\pi_1,\pi_2\}.
			\]
			Messages are symbolic terms
			\[
			t \in \mathcal{T}_{\Sigma}(X \cup \mathcal{N} \cup \mathcal{K}),
			\]
			where $X$ is a set of variables, $\mathcal{N}$ a set of nonces, and
			$\mathcal{K}$ a set of keys.
		};
		
		\node[block, below=4mm of sym1] (sym2) {%
			\textbf{Cryptographic operators as ideal constructors/destructors.}\\[2pt]
			Semantics is specified by an equational theory $E$ (or rewrite system), e.g.
			\[
			E \;=\;\bigl\{\,\mathsf{Dec}(\mathsf{Enc}(m,k),k)=m\,\bigr\}.
			\]
			Thus correctness is algebraic and exact, rather than probabilistic.
			Bit-level representation, entropy, cost, and implementation effects are
			abstracted away.
		};
		
		\node[block, below=4mm of sym2] (sym3) {%
			\textbf{Adversary as symbolic deduction.}\\[2pt]
			Given initial knowledge $\mathcal{K}_0$, the adversary derives new terms by
			closure under public operations and the equations in $E$:
			\[
			\mathcal{K}_{i+1} \;=\; \mathrm{cl}_{E}(\mathcal{K}_i),
			\qquad
			\mathcal{K}^{*} \;=\; \bigcup_{i\ge 0}\mathcal{K}_i.
			\]
			A ciphertext $\mathsf{Enc}(m,k)$ yields $m$ only when the adversary can
			derive the \emph{entire} key term $k$.
		};
		
		\node[block, below=4mm of sym3] (sym4) {%
			\textbf{Typical proof obligations.}\\[2pt]
			\[
			\text{Secrecy: } s \notin \mathcal{K}^{*}
			\qquad\qquad
			\text{Authentication: } \neg\,\mathsf{BadTrace}
			\]
			\[
			\text{Privacy / indistinguishability: } P \approx Q
			\]
			Hence symbolic methods are especially effective for protocol logic,
			message-flow analysis, trace properties, equivalence properties, and
			state-machine inconsistencies.
		};
		
		\node[smallblock, below=6mm of sym4] (sym5) {%
			\centering
			\textbf{Mathematical character}\\[2pt]
			Algebraic, term-based, proof-search friendly, highly automatable.
		};
		
		% =========================================================
		% RIGHT: COMPUTATIONAL
		% =========================================================
		\node[paneltitle] (comptitle) at (10.8,0) {Computational Security};
		
		\node[block, below=6mm of comptitle] (comp1) {%
			\textbf{Message universe as bitstrings.}\\[2pt]
			Messages live in
			\[
			m,c,k,r \in \{0,1\}^{*},
			\]
			and primitives are randomized polynomial-time algorithms operating on
			bitstrings with an explicit security parameter $\lambda$.
		};
		
		\node[block, below=4mm of comp1] (comp2) {%
			\textbf{Cryptographic primitive as an algorithmic tuple.}\\[2pt]
			A symmetric encryption scheme is
			\[
			\Pi=(\mathsf{Gen},\mathsf{Enc},\mathsf{Dec}),
			\]
			with
			\[
			k \leftarrow \mathsf{Gen}(1^\lambda), \qquad
			c \leftarrow \mathsf{Enc}_k(m;r), \qquad
			m' \leftarrow \mathsf{Dec}_k(c).
			\]
			Correctness is quantified over all legal keys and messages:
			\[
			\Pr\bigl[\mathsf{Dec}_k(\mathsf{Enc}_k(m;r)) = m\bigr]=1.
			\]
		};
		
		\node[block, below=4mm of comp2] (comp3) {%
			\textbf{Adversary as a PPT machine.}\\[2pt]
			The attacker is a probabilistic polynomial-time algorithm
			\[
			\mathcal{A} \in \mathrm{PPT}.
			\]
			Security is defined via experiments, reductions, and resource bounds.
			A canonical example is IND-CPA security:
			\[
			\mathrm{Adv}^{\mathrm{ind\mbox{-}cpa}}_{\Pi}(\mathcal{A},\lambda)
			=
			\left|
			\Pr\!\big[\mathsf{Exp}^{\mathrm{ind\mbox{-}cpa}}_{\Pi,\mathcal{A}}(\lambda)=1\big]
			-\tfrac12
			\right|.
			\]
		};
		
		\node[block, below=4mm of comp3] (comp4) {%
			\textbf{Typical proof style.}\\[2pt]
			One proves, for all $\mathcal{A}\in\mathrm{PPT}$,
			\[
			\mathrm{Adv}^{\mathrm{ind\mbox{-}cpa}}_{\Pi}(\mathcal{A},\lambda)
			\le
			\varepsilon(\lambda),
			\]
			or in concrete security,
			\[
			\forall\,\mathcal{A}\text{ with time }\le t:\quad
			\mathrm{Adv}_{\Pi}(\mathcal{A}) \le \varepsilon.
			\]
			Proofs proceed by game hopping, reductions to hardness assumptions, and
			explicit accounting of probability and running time.
		};
		
		\node[smallblock, below=6mm of comp4] (comp5) {%
			\centering
			\textbf{Mathematical character}\\[2pt]
			Probabilistic, complexity-theoretic, reduction-based, closer to
			mainstream provable cryptography.
		};
		
		% =========================================================
		% FLOWS
		% =========================================================
		\draw[arrow] (symtitle) -- (sym1);
		\draw[arrow] (sym1) -- (sym2);
		\draw[arrow] (sym2) -- (sym3);
		\draw[arrow] (sym3) -- (sym4);
		\draw[arrow] (sym4) -- (sym5);
		
		\draw[arrow] (comptitle) -- (comp1);
		\draw[arrow] (comp1) -- (comp2);
		\draw[arrow] (comp2) -- (comp3);
		\draw[arrow] (comp3) -- (comp4);
		\draw[arrow] (comp4) -- (comp5);
		
		% =========================================================
		% COMPARISON / BOTTOM
		% =========================================================
		\node[block, text width=15.7cm, align=center, below=11mm of $(sym5)!0.5!(comp5)$] (trade) {%
			\textbf{Conceptual trade-off}\\[3pt]
			\emph{Symbolic security} replaces cryptographic computation by algebraic structure
			$(\mathcal{T}_{\Sigma},E)$, gaining tractability and automation. \qquad
			\emph{Computational security} models actual randomized algorithms on bitstrings,
			gaining fidelity to standard provable-security notions, but usually at a higher
			proof and automation cost.
		};
		
		\draw[compare] (sym5.east) -- node[above, align=center] {\textbf{abstraction $\;\longleftrightarrow\;$ fidelity}} (comp5.west);
		
		% =========================================================
		% BACKGROUND GROUPS
		% =========================================================
		\begin{scope}[on background layer]
			\node[draw, rounded corners=6pt, line width=0.9pt,
			fit=(symtitle)(sym5), inner sep=8pt] {};
			\node[draw, rounded corners=6pt, line width=0.9pt,
			fit=(comptitle)(comp5), inner sep=8pt] {};
		\end{scope}
		
	\end{tikzpicture}
	
\end{document}